\chapter{系统功能模块划分}

\section{功能需求分析}

根据需求分析，系统可以划分为以下六个主要功能模块：

\begin{table}[H]
  \centering
  \caption{系统功能模块划分}
  \label{tab:modules}
  \begin{tabularx}{\textwidth}{|l|X|}
    \hline
    \textbf{模块代码} & \textbf{功能需求} \\
    \hline
    M1 基础数据管理 & 门店、桌游目录、桌位主数据及 URL 维护 \\
    \hline
    M2 桌位平面图 & 空闲/预约中/占用状态可视化展示 \\
    \hline
    M3 预约管理 & 新建预约、冲突检测、取消、列表查询 \\
    \hline
    M4 对局与计费 & 入场开台、现场开台、结算关台 \\
    \hline
    M5 战绩与排行 & 战绩写入、排行榜维护、聚合统计 \\
    \hline
    M6 统计报表 & 存储过程支撑收入、利用率等分析 \\
    \hline
  \end{tabularx}
\end{table}

\section{数据流图（DFD）分析}

\subsection{0 层数据流图（上下文图）}

0 层数据流图展示系统与外部实体的交互关系。

\begin{figure}[H]
  \centering
  \begin{tikzpicture}[scale=1.2]
    % 系统
    \draw[fill=blue!30, draw=blue, thick] (2,1.5) rectangle (4,3.5) node[midway, align=center] {\textbf{桌游预约与\\战绩系统}};

    % 外部实体
    \draw[fill=green!30, draw=green, thick] (0.5,4) circle (0.4) node[midway] {\small 顾客};
    \draw[fill=green!30, draw=green, thick] (0.5,2.5) circle (0.4) node[midway] {\small 会员};
    \draw[fill=green!30, draw=green, thick] (0.5,1) circle (0.4) node[midway] {\small 管理员};
    \draw[fill=yellow!20, draw=orange, thick] (5.5,2.5) rectangle (6.5,3.5) node[midway] {\small MySQL};

    % 数据流
    \draw[->, thick] (0.9,4) -- (2,3.2) node[midway, above] {\small 预约查询};
    \draw[->, thick] (0.9,2.5) -- (2,2.5) node[midway, above] {\small 预约/战绩};
    \draw[->, thick] (0.9,1) -- (2,1.8) node[midway, above] {\small 维护报表};
    \draw[->, thick] (4,3.2) -- (5.5,3) node[midway, above] {\small 持久化};
    \draw[->, thick] (5.5,2.7) -- (4,2.3) node[midway, below] {\small 查询};
    \draw[->, thick] (4,1.8) -- (2,1.5) node[midway, below] {\small 反馈};
  \end{tikzpicture}
  \caption{0 层数据流图（上下文图）}
  \label{fig:dfd0}
\end{figure}

\subsection{1 层数据流图（功能分解）}

1 层数据流图展示系统内部的主要加工与数据存储。

\begin{figure}[H]
  \centering
  \begin{tikzpicture}[scale=0.9]
    % 外部实体
    \draw[fill=green!30, draw=green, thick] (0.5,6) circle (0.35) node[midway, font=\small] {顾客};
    \draw[fill=green!30, draw=green, thick] (0.5,3) circle (0.35) node[midway, font=\small] {会员};
    \draw[fill=green!30, draw=green, thick] (0.5,0) circle (0.35) node[midway, font=\small] {管理员};

    % 加工
    \draw[fill=blue!30, draw=blue, thick] (2.5,6) rectangle (3.5,6.8) node[midway, font=\small] {1.0 预约};
    \draw[fill=blue!30, draw=blue, thick] (2.5,4.5) rectangle (3.5,5.3) node[midway, font=\small] {2.0 开台};
    \draw[fill=blue!30, draw=blue, thick] (2.5,3) rectangle (3.5,3.8) node[midway, font=\small] {3.0 结算};
    \draw[fill=blue!30, draw=blue, thick] (2.5,1.5) rectangle (3.5,2.3) node[midway, font=\small] {4.0 战绩};
    \draw[fill=blue!30, draw=blue, thick] (2.5,0) rectangle (3.5,0.8) node[midway, font=\small] {5.0 统计};

    % 数据存储
    \draw[fill=yellow!20, draw=orange, thick] (5.5,6) rectangle (6.5,6.8) node[midway, font=\small] {D1 预约};
    \draw[fill=yellow!20, draw=orange, thick] (5.5,4.5) rectangle (6.5,5.3) node[midway, font=\small] {D2 对局};
    \draw[fill=yellow!20, draw=orange, thick] (5.5,3) rectangle (6.5,3.8) node[midway, font=\small] {D3 战绩};
    \draw[fill=yellow!20, draw=orange, thick] (5.5,1.5) rectangle (6.5,2.3) node[midway, font=\small] {D4 排行};

    % 数据流
    \draw[->, thick] (0.85,6) -- (2.5,6.4);
    \draw[->, thick] (0.85,3) -- (2.5,4.9);
    \draw[->, thick] (0.85,0) -- (2.5,0.4);
    \draw[<->, thick] (3.5,6.4) -- (5.5,6.4);
    \draw[<->, thick] (3.5,4.9) -- (5.5,4.9);
    \draw[<->, thick] (3.5,3.4) -- (5.5,3.4);
    \draw[<->, thick] (3.5,1.9) -- (5.5,1.9);
  \end{tikzpicture}
  \caption{1 层数据流图（功能分解）}
  \label{fig:dfd1}
\end{figure}

\section{软件功能模块图}

系统采用自顶向下的模块划分方式，如图 \ref{fig:module_tree} 所示。

\begin{figure}[H]
  \centering
  \begin{tikzpicture}[
    level distance=1.5cm,
    sibling distance=2.5cm,
    every node/.style={draw, rounded rectangle, minimum width=2cm, minimum height=0.6cm, font=\small},
    level 1/.style={sibling distance=3.5cm},
    level 2/.style={sibling distance=2cm}
  ]
    \node[fill=blue!30] {桌游运营系统}
      child {node[fill=cyan!30] {M1 基础数据}
        child {node[fill=cyan!20] {门店/目录/桌位}}
      }
      child {node[fill=cyan!30] {M2 平面图}
        child {node[fill=cyan!20] {状态着色}}
      }
      child {node[fill=cyan!30] {M3 预约}
        child {node[fill=cyan!20] {新建/取消/列表}}
      }
      child {node[fill=cyan!30] {M4 对局计费}
        child {node[fill=cyan!20] {入场/开台/结算}}
      }
      child {node[fill=cyan!30] {M5 战绩排行}
        child {node[fill=cyan!20] {录入/排行榜}}
      }
      child {node[fill=cyan!30] {M6 统计报表}
        child {node[fill=cyan!20] {存储过程}}
      };
  \end{tikzpicture}
  \caption{软件功能模块树}
  \label{fig:module_tree}
\end{figure}

\section{模块间的数据依赖关系}

\begin{table}[H]
  \centering
  \caption{模块间数据依赖关系}
  \label{tab:module_deps}
  \begin{tabularx}{\textwidth}{|l|X|}
    \hline
    \textbf{模块} & \textbf{依赖关系} \\
    \hline
    M1 基础数据 & 无依赖（基础模块） \\
    \hline
    M2 平面图 & 依赖 M1（需要桌位信息） \\
    \hline
    M3 预约 & 依赖 M1、M2（需要桌位和状态） \\
    \hline
    M4 对局计费 & 依赖 M1、M3（需要桌位和预约） \\
    \hline
    M5 战绩排行 & 依赖 M1、M4（需要游戏和对局） \\
    \hline
    M6 统计报表 & 依赖 M1、M3、M4、M5（综合查询） \\
    \hline
  \end{tabularx}
\end{table}
